\chapter{Methodology}
\label{ch:methodology}

This chapter describes the experimental design (Section~\ref{sec:m-design}) and the metric design
(Sections~\ref{sec:m-lens} to~\ref{sec:m-metaeval}). The guiding principle is that every quantity is
anchored to ground truth. Because the Hubble suite ships, for each fact, a model trained to memorise it
(the knows reference) and a model trained without it (the never-knew reference), we can calibrate both the
internal readout and the metric derived from it against controlled knowledge states rather than against
guesswork.

\section{Experimental Design}
\label{sec:m-design}

\subsection{The Hubble suite and the six models}
\label{sec:m-models}

All experiments use the Hubble benchmark \citep{allegro2025hubble}, an open suite of Llama-3 style models
pretrained with a controlled set of synthetic facts inserted. We work at two scales, 1B (16 layers,
hidden size 2048) and 8B (36 layers, hidden size 4096), and with two reference models per scale. The
\emph{Standard} model, trained without the inserted facts (a never-knew lower bound), and the
\emph{Perturbed} model, trained to memorise them (a knows upper bound). To these we add four unlearned
models, one per method family of Section~\ref{sec:method-families}, namely NPO \citep{zhang2024npo}, GradDiff
\citep{jang2022knowledge}, IDK \citep{maini2024tofu}, and RMU \citep{li2024wmdp}. RMU follows the Hubble
configuration (intervention layers $[5,6,7]$ at 8B, $[3,4,5]$ at 1B); the other three are extensions not
present in the original paper. Each method is tuned so that unlearning is \emph{selective}: the target is
forgotten while unrelated knowledge is preserved.

The planted facts are synthetic biographies of fictional people with attributes such as birthdate and
nationality. Following Hubble, the data are split into an \emph{Unlearn} (forget) set, a \emph{Keep}
(retain) set of in-domain near-neighbours, and an \emph{Unseen} held-out set. Crucially, we use the Keep
biographies as the retain set during unlearning (Hubble's Appendix-G ablation) rather than generic web
text; this is the deliberate choice that makes forgetting \emph{separable} from retention for the internal
read. Unless stated otherwise, the study targets the birthdate attribute, which is unguessable and cleanly
tokenised, over $n = 44$ forget-set facts.

\paragraph{Data reuse.}
We state the data flow plainly. The same $44$ forget-set facts are used throughout. They are the facts the
models are unlearned on, the facts a small grid search tunes each method to forget while preserving the
Keep set (Appendix~\ref{app:hp}), and the facts on which the output metrics, $D_\text{mech}$, the paraphrase
analysis, and the robustness attacks are all computed. There is no held-out set of \emph{new} forget facts.
The results should therefore be read as an analysis of these trained benchmark models, that is, of where
each method places the target knowledge internally and how the metrics score it, rather than as an
out-of-sample estimate of performance on unseen facts. The leave-one-fact-out anchor construction
(Section~\ref{sec:m-metric}) removes only self-comparison within this fixed set, not this reuse.

\subsection{The probing task}
\label{sec:m-probe}

We query a fact with an \emph{infill} probe. The model is given a biography that stops immediately before
the sensitive attribute, for example \emph{``\ldots{}Billy Fisher was born on \_\_\_''}, and is asked to
supply the continuation. Following Hubble's own loss-based-choice protocol, infill accuracy scores the
true birthdate against nine distractors and counts the fact recalled when the true answer is the
highest-likelihood option (chance $\approx 0.10$). Four prompt formats provide more or less surrounding
context (name-only, intro-prefix, full-prefix, and full-prefix with a suffix); the internal analysis uses
the leak-free full-prefix format so that the model has the context it needs to recall the fact.

\subsection{Reading the internal computation with three lenses}
\label{sec:m-lens}

To read what the model represents internally we decode intermediate hidden states into vocabulary space
with a lens (Section~\ref{sec:mech-interp}) and record the \emph{rank} of the gold answer token at each
layer, from rank~1 (top prediction) to rank in the tens of thousands (buried). Aggregating over the gold
sub-token positions and taking the median over facts yields a rank-by-layer trajectory
$r = (r_0, r_1, \ldots, r_L)$.

\paragraph{The logit lens and two corrections.}
The primary readout is the logit lens \citep{nostalgebraist2020logitlens}: the model's own final
normalisation and unembedding applied to each intermediate state. Two implementation details proved
essential for correctness, and all reported numbers use the corrected readout. First,
\emph{double-normalisation}: a transformer's last hidden state is already layer-normalised, so re-applying
the norm at the final layer distorts it and produces a spurious final-layer rank spike; we instead use the
model's true output logits at the last layer. Second, \emph{boundary-safe tokenisation}: because of
byte-level tokenisation, a trailing space in the prompt causes the na\"ive code to track the wrong first
answer token (\texttt{"July"} rather than the space-prefixed \texttt{" July"} the model actually predicts),
inflating ranks; we strip and re-join the prompt so that the tracked token is the one the model produces.
Together these fixes reduce mid-stack ranks by roughly five- to forty-fold and remove the false end-spike
(for example, IDK's 8B forget rank falls from 958 to 9).

\paragraph{The tuned lens.}
For comparison we also train a tuned lens \citep{belrose2023tunedlens}: one affine map per layer,
identity-initialised, optimised on held-out text so that the readout reproduces the model's final
distribution. We train to a per-layer Kullback-Leibler divergence of roughly $0.05$ on every model. The
tuned lens is used only as a cross-check, for reasons made precise in Section~\ref{sec:m-whylogit} and
demonstrated empirically in Chapter~\ref{ch:results}.

\paragraph{The Jacobian lens.}
Finally we implement the Jacobian lens \citep{anthropic2026workspace} as a first-order,
interventionally-grounded readout. It transports a
hidden state through the model's own averaged input-output Jacobian, taken over $24$ reference contexts
sampled from a pool of $38$ generic sentences (leave-the-query-out), before decoding it with the unembedding. We compute the required Jacobian-vector products with a
hook-injection scheme (injecting the state into the model's own forward pass) and a double-backward
formulation, verifying on every layer that the injection reproduces the true forward pass exactly. Two
choices are needed for a usable readout. Restoring the first-order Taylor \emph{intercept} (without it the
trajectory is jagged), and averaging the Jacobian over \emph{diverse} reference text rather than
same-template prompts (without which the readout trivially reconstructs the model's own output at every
layer). The Jacobian lens is expensive and serves as a confirmatory, interventionally-grounded view rather
than the basis of the metric.

\section{Designing the Metric \texorpdfstring{$D_\text{mech}$}{D\_mech}}
\label{sec:m-metric}

From the logit-lens trajectory we derive a scalar metric of internal knowledge. For a fact with trajectory
$r$, define two anchor trajectories from the ground-truth reference models: $K$, the knows anchor
(Perturbed), and $\varnothing$, the never-knew anchor (Standard). Over the mid-to-late band $B$ (layers
$18$ to $36$ at 8B, the top half of the network, where the logit lens is reliable; fixed a priori as the top
half rather than tuned on the outcome, with its sensitivity to the band choice reported in
Section~\ref{sec:r-crossbench}), the \emph{dual-anchor mechanistic distance} is
\begin{equation}
  D_\text{mech} \;=\; \frac{1}{|B|}\sum_{L\in B}\bigl|\log r_L - \log \varnothing_L\bigr|
                    \;-\; \frac{1}{|B|}\sum_{L\in B}\bigl|\log r_L - \log K_L\bigr|.
  \label{eq:dmech}
\end{equation}
A positive $D_\text{mech}$ means the trajectory lies closer to the knows template than to the never-knew
template within the selected band, whatever the output says; we read it as an indicator of residual internal
accessibility, not as independent proof of stored knowledge. Three design choices define the metric.

\paragraph{Why the logarithm.}
Ranks live on a multiplicative scale, spanning five orders of magnitude. A raw difference $|5000-4000|$
would dominate $|1-10|$, yet rank~1 versus rank~10 is hugely meaningful (top guess versus tenth) while
rank~4000 versus~5000 is not (both hopelessly buried). Taking logarithms turns the distance into a
rank-\emph{ratio}, so that improving from rank~10 to rank~1 counts far more than from rank~5000 to~4000,
and makes the per-layer terms comparable so that the band average is not dominated by one buried layer.

\paragraph{Why two anchors.}
$D_\text{mech}$ is a \emph{relative} quantity. It asks not ``how buried is the fact?'' but ``is the
trajectory closer to knows or to never-knew?''. This relative construction is the source of the metric's
improved paraphrase stability (Section~\ref{sec:m-metaeval}). A meaning-preserving rewording nudges the whole
trajectory a little, changing its distance to \emph{both} anchors by approximately the same amount, and the
subtraction cancels most of that common-mode shift while preserving the knowledge signal. As
Chapter~\ref{ch:results} shows, a na\"ive single-anchor internal metric (the minimum rank over the band)
lacks this cancellation and is substantially more sensitive to paraphrasing.

\paragraph{Leave-one-fact-out anchors.}
To avoid circularity, the anchors for a given fact are the median trajectories of the \emph{other} $43$
facts (leave-one-fact-out), so a fact is never compared against itself. The tight distributions of the Full
and Retain cells (Section~\ref{sec:m-metaeval}) confirm that these medians are representative templates.

\subsection{Why the logit lens grounds the metric}
\label{sec:m-whylogit}

Of the three lenses, only the logit lens is a parameter-free, exactly reproducible function of the model:
it introduces no trained or estimated components, so $D_\text{mech}$ is a property of the model rather than
of our modelling choices. The tuned lens, being trained to reproduce the \emph{output}, tends to surface
the answer at every layer and therefore collapses the very trajectory differences $D_\text{mech}$ reads
(Chapter~\ref{ch:results}); the Jacobian lens introduces a dependence on the choice of reference corpus.
Moreover, $D_\text{mech}$ only uses the band where the logit lens is reliable, side-stepping its sole
weakness at early layers. We therefore build the metric on the logit lens and reserve the tuned and
Jacobian lenses for cross-checking.

\section{The Faithfulness Meta-Evaluation}
\label{sec:m-metaeval}

Whether $D_\text{mech}$ is \emph{better} than existing metrics is a question one level up. Which metric is
more \emph{faithful}, that is, tracks knowledge rather than surface behaviour? We answer it with the
faithful analysis based on OpenUnlearning \citep{dorna2025openunlearning}, which judges a metric on two
criteria that must both hold: \emph{Knowledge Discrimination} (the metric must score models that hold the
target knowledge differently from models that lack it) and \emph{Semantic Invariance} (the metric must give
the same score when the question is reworded but the knowledge is unchanged). Discrimination is necessary
but easy; invariance is the hard criterion on which lexical metrics fail.

\paragraph{The competing metrics.}
We pit $D_\text{mech}$, and the na\"ive internal metric $\min$-rank, against nine established output and
lexical metrics computed per fact. Gold-answer probability, ROUGE-L, exact-match memorisation,
extraction strength, truth ratio, and four membership-inference scores (loss, Min-$K$\%, Min-$K$\%++, and
zlib-normalised loss).

\paragraph{The four pools.}
Following the meta-evaluation construction, we build four cells on the Hubble suite, each a model evaluated
on the birthdate of every entity ($n=44$):
\begin{itemize}
  \item \textbf{Full} (Positive): Perturbed model, canonical prompt, correct answer: \emph{knows}.
  \item \textbf{Paraphrased} (Positive): Perturbed model, reworded prompt, correct answer: \emph{knows}
  the same fact in different words. The rewording changes only the trigger clause (``\ldots{}was born on''
  becomes ``\ldots{}has a recorded birth date of''), keeping all surrounding context.
  \item \textbf{Retain} (Negative): Standard model, correct answer, knowledge \emph{absent} (never knew).
  Here ``Retain'' denotes this never-knew reference pool; it is distinct from the \emph{Keep}/retain-utility
  facts used elsewhere to check collateral damage, which are preserved neighbour facts rather than a
  never-knew reference.
  \item \textbf{Perturbed-wrong} (Negative): Perturbed model scored against a \emph{wrong} birthdate:
  the model holds no such fact, catching metrics that fire on fluency or format rather than the specific
  fact.
\end{itemize}
The positive pool $P = \{\text{Full}, \text{Paraphrased}\}$ and negative pool $N = \{\text{Retain},
\text{Perturbed-wrong}\}$ each contain $88$ instances. Comparing Full with Paraphrased isolates invariance
(same model, same knowledge, only the words change); comparing $P$ with $N$ measures discrimination.

\paragraph{Statistics and assumptions.}
With $n=44$ and bounded, skewed, heavy-tailed metric distributions, we use non-parametric tests
throughout. A within-$P$ paired Wilcoxon signed-rank test, a between-pool Mann-Whitney $U$ test, and
Cliff's delta effect sizes; we also report a TOST equivalence test for completeness, with an exploratory,
non-pre-registered equivalence bound that no claim rests on. The effect sizes are the headline numbers:
$\text{invar\_shift} = |\delta(\text{Full},\text{Paraphrased})|$ (0 = no shift) and $\text{discrim} =
|\delta(P,N)|$ (1 = perfect separation), combined into a convenience score $\text{faithful} = \text{discrim}
\cdot (1 - \text{invar\_shift})$. We call a metric \emph{faithful} when it is both discriminative and
paraphrase-stable, which we operationalise as a discrimination effect size above $0.5$ together with an
invariance shift below $0.5$; on Hubble-8B only $D_\text{mech}$ meets both thresholds. A methodological
caveat we treat carefully. At $n=44$ the within-$P$
Wilcoxon $p$-value \emph{saturates}, flagging even a tiny consistent shift as significant for \emph{every}
metric, including $D_\text{mech}$; we therefore judge invariance by \emph{effect size}, not by a
$p<0.05$ threshold, and report the $p$-values only for completeness. The facts are treated as the unit of
analysis, but they share model checkpoints, prompt templates, and overlapping leave-one-out anchors, so the
tests are benchmark-conditional rather than estimates over an unrestricted population of facts. As a further
external-validity limit we state honestly, the paraphrase and the wrong-target are each a single fixed
realisation rather than sampled across styles; a multi-paraphrase robustness check
(Section~\ref{sec:r-invariance}) partially addresses this.
